\documentclass[SCMIJ,viewpoint ,submit,moreauthors,pdftex]{Files/scmij} 
\firstpage{1} 
\makeatletter 
\setcounter{page}{\@firstpage} 
\makeatother




\Title{Title of the article for SCMI journal submission }

% Author Orchid ID: enter ID or remove command
\newcommand{\orcidauthorA}{0000-0000-000-000X} % Add \orcidA{} behind the author's name
\newcommand{\orcidauthorB}{0000-0000-000-000X} % Add \orcidB{} behind the author's name

% Authors, for the paper (add full first names)
\Author{First Author $^{a}$\orcidA{}, Second Author $^{a}$\orcidB{} and Last Author $^{b,}$*}

% Affiliations / Addresses (Add [a] after \address if there is only one affiliation.)
\address{%
$^{(a)}$ \quad Full address of first Affiliation (a) including city and country\\
$^{(b)}$ \quad Full address of first Affiliation (b) including city and country}

% Contact information of the corresponding author
\corres{Corresponding author: Email address of the Corresponding author}

%**********************************%
\abstract{The abstract is a single paragraph of about 200 words maximum. For research articles, abstracts should give a pertinent overview of the work. We strongly encourage authors to use the following style of structured abstracts, without headings: Background, Methods, Results, and Conclusion. The abstract should be an objective representation of the article; it must not contain results that are not presented and substantiated in the main text and should not exaggerate the main conclusions. The abstract is supposed to cover background and existed issues (10\% of abstract), what you have proposed (15\%), describe the solution procedure in detail (45\%), and describe the experiment's performance, efficiency, etc. (30\%).  }


\keyword{keyword 1; keyword 2; keyword 3; keyword 3 }

%**********************************%

\begin{document}

\section{Introduction}
The introduction should concisely place the study in a comprehensive meaning and highlight why it is essential. It should describe the scope of the work and its importance. The current state of the research field should be evaluated thoroughly, and key publications cited. Please highlight controversial and diverging hypotheses when necessary. Conclusively state the main aim of the work and highlight the principal conclusions. As far as possible, please keep the introduction comprehensible to scientists outside your particular field of research. Citing a journal paper \cite{ref-journal}. Moreover, now citing a book reference \cite{ref-book}. 


%%%%%%%%%%%%%%%%%%%%%%%%%%%%%%%%%%%%%%%%%%
\section{Results}

This section may be divided by subheadings. It should provide a concise and precise description of the experimental results, their interpretation, and the experimental conclusions are drawn.
\begin{quote}
This section may be divided by subheadings. It should provide a concise and precise description of the experimental results, their interpretation, and the experimental conclusions are drawn.
\end{quote}

%%%%%%%%%%%%%%%%%%%%%%%%%%%%%%%%%%%%%%%%%%
\subsection{Subsection}
\unskip
\subsubsection{Subsubsection}

Bulleted lists look like this:
\begin{itemize}[leftmargin=*,labelsep=5.8mm]
\item	First bullet
\item	Second bullet
\item	Third bullet
\end{itemize}

Numbered lists can be added as follows:
\begin{enumerate}[leftmargin=*,labelsep=4.9mm]
\item	First item 
\item	Second item
\item	Third item
\end{enumerate}

The text continues here.

\subsection{Figures, Tables, and Schemes}

All figures and tables should be cited in the main text as Figure 1, Table 1, etc.

\begin{figure}[H]
\centering
\includegraphics[width=8 cm]{Figs/scmij.PNG}
\caption{This is a figure; schemes follow the same formatting. If there are multiple panels, they should be listed as (\textbf{a}) Description of what is contained in the first panel. (\textbf{b}) Description of what is contained in the second panel.}
\end{figure}   
 
 Figures should be placed in the main text near to the first time they are cited. A caption on a single line should be centered. Text Text Text Text Text Text Text Text Text Text Text Text



\begin{table}[H]
\caption{This is a table caption. }
\centering
%% \tablesize{} %% You can specify the fontsize here, e.g., \tablesize{\footnotesize}. If commented out \small will be used.
\begin{tabular}{ccc}
\toprule
\textbf{Title 1}	& \textbf{Title 2}	& \textbf{Title 3}\\
\midrule
entry 1		& data			& data\\
entry 2		& data			& data\\
\bottomrule
\end{tabular}
\end{table}

Tables should be placed in the main text near to the first time they are cited.  Text Text Text Text Text Text Text Text Text Text Text Text
Text

%\begin{listing}[H]
%\caption{Title of the listing}
%\rule{\textwidth}{1pt}
%\raggedright Text of the listing. In font size footnotesize, small, or normalsize. Preferred format: left aligned and single spaced. Preferred border format: top border line and bottom border line.
%\rule{\textwidth}{1pt}
%\end{listing}


\subsection{Formatting of Mathematical Components}

This is an example of an equation:

\begin{equation}
x \times y \sum z = d
\end{equation}

Please punctuate equations as regular text. Theorem-type environments (including propositions, lemmas, corollaries, etc.) can be formatted as follows:
%% Example of a theorem:
\begin{Theorem}
Example text of a theorem.
\end{Theorem}

The text continues here. Proofs must be formatted as follows:

%% Example of proof:
\begin{proof}[Proof of Theorem 1]
Text of the proof. Note that the phrase `of Theorem 1' is optional if it is clear which theorem is being referred to.
\end{proof}
The text continues here.


\section{Discussion}

Writers should present the results and how they can be interpreted from former studies and effective hypotheses. The conclusions and their associations should be discussed in the most comprehensive context possible. Future research directions may also be highlighted.


\section{Conclusions}

The conclusion is a summary of the main ideas that come out of the discussion. This section is not mandatory but can be added to the manuscript if the discussion is unusually long or complex.


%%%%%%%%%%%%%%%%%%%%%%%%%%%%%%%%%%%%%%%%%%
\vspace{6pt} 


%****************************************
\acknowledgments{In this section, the authors can acknowledge any support given. This may include administrative and technical support or funding. .}


%****************************************

%% optional
\appendixtitles{no} 
\appendix
\section{}
\unskip
\subsection{}
The appendix is an elective segment that can include features and data supplemental to the main text. For example, descriptions of experimental features that would disrupt the main text's flow remain crucial to understanding and reproducing the research shown; figures of replicates for experiments of which representative data is shown in the main text can be added here brief, or as Supplementary data. Mathematical proofs of results not central to the paper can be added as an appendix.

\section{}
All appendix sections must be cited in the main text. In the appendixes, Figures, Tables, etc., should be labeled starting with `A,' e.g., Figure A1, Figure A2, etc. 

%****************************************
%****************************************


\reftitle{References}


%=====================================
% References, variant A: external bibliography
%=====================================
%\externalbibliography{yes}
%\bibliography{your_external_BibTeX_file}

%=====================================
% References, variant B: internal bibliography
%=====================================
\begin{thebibliography}{999}
% Reference 1
\bibitem[Author1(year)]{ref-journal}
Author1, T. The title of the cited article. {\em Journal Abbreviation} {\bf 2008}, {\em 10}, 142--149.
\bibitem[Author2(year)]{ref-book}
Author2, L. The title of the cited contribution. In {\em The Book Title}; Editor1, F., Editor2, A., Eds.; Publishing House: City, Country, 2007; pp. 32--58.
\bibitem[Author2(year)]{ref-book}
Author2, L. The title of the cited contribution. In {\em The Book Title}; Editor1, F., Editor2, A., Eds.; Publishing House: City, Country, 2007; pp. 32--58.
% Reference 2
\bibitem[Author2(year)]{ref-book}
Author2, L. The title of the cited contribution. In {\em The Book Title}; Editor1, F., Editor2, A., Eds.; Publishing House: City, Country, 2007; pp. 32--58.
\end{thebibliography}


%%%%%%%%%%%%%%%%%%%%%%%%%%%%%%%%%%%%%%%%%%
\end{document}
